\chapter{Ergebnisse}
\label{ch:ergebnisse}

Die folgenden Kennzahlen wurden aus dem aktuellen Datensatz und dem bestehenden Code abgeleitet. Alle Streckenkennwerte bleiben zwischen den beiden Glättungslaeufen gleich, weil sie nur von den GPS-Daten abhängen. Unterschiede entstehen ausschliesslich in Geschwindigkeit, Beschleunigung und den daraus folgenden elektrischen Groessen.

\section{Streckenkennwerte}

\begin{table}[htbp]
\centering
\caption{Ermittelte Kennzahlen des aktuellen Datensatzes}
\label{tab:streckenkennzahlen}
\begin{tabular}{@{}lr@{}}
\toprule
Kennwert & Wert \\
\midrule
GPS-Punkte & 2284 \\
Gesamtdauer & 16372.391\,s \\
Gesamtstrecke & 94.267\,km \\
Hoehenmeter aufwaerts & 1095.895\,m \\
Hoehenmeter abwaerts & 1097.146\,m \\
Durchschnittsgeschwindigkeit & 20.728\,km/h \\
Maximale Rohgeschwindigkeit & 48.772\,km/h \\
Maximale geglaettete Geschwindigkeit & 48.584\,km/h \\
Maximale Rohbeschleunigung & 2.982\,m/s\textsuperscript{2} \\
Minimale Rohbeschleunigung & -3.790\,m/s\textsuperscript{2} \\
\bottomrule
\end{tabular}
\end{table}

\section{Vergleich der Glättung}

\begin{table}[htbp]
\centering
\caption{Vergleich ohne und mit Glättung}
\label{tab:vergleich-glaettung}
\begin{tabularx}{\textwidth}{@{}lXX@{}}
\toprule
Kennwert & Glättung aus & Glättung an \\
\midrule
Aktive Geschwindigkeitsspitze & 48.772\,km/h & 48.584\,km/h \\
Aktive Beschleunigungsspitze & 2.982\,m/s\textsuperscript{2} & 0.583\,m/s\textsuperscript{2} \\
Maximale mechanische Leistung & 2515.913\,W & 974.146\,W \\
Maximaler Motorstrom & 61.628\,A & 24.843\,A \\
End-SoC LiPo & 0.0\,\% & 0.0\,\% \\
End-SoC NMC & 0.0\,\% & 0.0\,\% \\
Minimale Spannung LiPo & 31.278\,V & 31.292\,V \\
Minimale Spannung NMC & 31.115\,V & 31.097\,V \\
Bremswiderstandsenergie & 0.000\,Wh & 0.000\,Wh \\
\bottomrule
\end{tabularx}
\end{table}

Die Strecke selbst bleibt unveraendert, aber die Glättung reduziert die Spitzen deutlich. Besonders der Unterschied bei Beschleunigung, Leistung und Motorstrom ist relevant: Ohne Glättung fuehrt ein kurzer Geschwindigkeitssprung zu einer sehr hohen Lastspitze. Mit Glättung wird der Verlauf deutlich ruhiger, waehrend die Strecke und die mittlere Geschwindigkeit praktisch gleich bleiben.

\section{Interpretation}

Beide Akkus werden im aktuellen Datensatz vollstaendig entladen. Das bedeutet nicht, dass die Akkus physikalisch identisch reagieren; vielmehr ist der Lastbedarf so hoch, dass beide Modelle vor dem Streckenende auf 0\,\% SoC fallen. Die NMC-Kurve liegt bei gleichem SoC typischerweise niedriger in der Spannung als die LiPo-Kurve, weshalb die minimale Spannung im NMC-Lauf niedriger ausfaellt.

Der wesentliche Nutzen der Glättung zeigt sich in den Spitzenwerten. Sie senkt die maximal benoetigte Leistung auf rund \SI{974}{\watt} statt \SI{2516}{\watt} und reduziert den maximalen Motorstrom auf etwa \SI{24.8}{\ampere} statt \SI{61.6}{\ampere}. Damit wird die numerische Auswirkung von GPS-Rauschen direkt sichtbar.
